\documentclass[UTF8]{ctexart}
\usepackage{geometry}
\geometry{a4paper, margin=2.5cm}
\usepackage{listings}
\usepackage{xcolor}
\usepackage{fontspec}
\usepackage{amsmath}
\usepackage{graphicx}
\usepackage{silence}
\usepackage{lmodern}
\WarningFilter{latex}{Font shape}
\WarningFilter{LaTeX}{font}
\setmonofont{DejaVu Sans Mono}
\lstset{
    basicstyle=\ttfamily\small,
    breaklines=true,
    frame=single,
    numbers=left,
    numberstyle=\tiny,
    tabsize=4,
    language=C,
    keywordstyle=\color{blue},
    commentstyle=\color{green!60!black},
    stringstyle=\color{red}
}

\title{微服务拓扑故障影响范围分析实验报告}
\author{}

\begin{document}
\maketitle

\begin{center}
    \textbf{班级：}2024217802 \quad \textbf{姓名：}严子竣 \quad \textbf{学号：}2024212622
\end{center}

\section{需求分析}

在分布式微服务架构中，服务之间的调用关系可以抽象为有向图，其中顶点代表微服务节点，有向边代表调用关系。当某个服务节点发生故障时，下游无法正常访问该节点的服务数量即为“故障影响范围”。本实验将该工程需求转化为有向图的可达性问题：随机生成符合微服务调用特征的有向图，计算从每个顶点出发所能到达的顶点数量的平均值与最大值，以此评估系统的故障扩散风险，为架构设计和容灾规划提供定量依据。

\subsection{程序功能}
\begin{itemize}
    \item 从终端交互接收两个正整数 \(V\)（顶点数）和 \(E\)（边数），取值范围分别为 \(V \geq 1\)，\(0 \leq E \leq V \times (V-1)\)。
    \item 根据微服务调用特征随机生成简单有向图：前 \(10\%\) 的顶点视为核心服务，其作为边的源顶点的概率更高（模拟核心服务被更多下游服务调用）。
    \item 对每个顶点执行广度优先搜索（BFS），计算从该顶点出发可到达的顶点数量（包括自身）。
    \item 统计所有顶点可达数的平均值和最大值，并依据阈值（平均值/最大值占比 \(\geq 50\%\)）给出故障扩散风险评估结论。
\end{itemize}

\subsection{输入形式与值的范围}
\begin{itemize}
    \item 程序运行后，会提示用户输入两个整数：
        \begin{itemize}
            \item 顶点数 \(V\)：正整数，取值范围为 \(1 \sim\) 任意正整数。
            \item 边数 \(E\)：非负整数，最大可能边数为 \(V \times (V-1)\)。
        \end{itemize}
    \item 若输入不合法（如 \(V \leq 0\)、\(E < 0\) 或 \(E > V(V-1)\)），程序会输出错误信息并退出。
\end{itemize}

\subsection{输出形式}
\begin{itemize}
    \item 输出所有顶点可达数的平均值（保留两位小数）、最大值。
    \item 根据阈值判断系统耦合性：
        \begin{itemize}
            \item 若平均值占比 \(\geq 50\%\) 或最大值占比 \(\geq 50\%\)，输出“说明系统耦合度高”；
            \item 否则输出“说明系统耦合度低”。
        \end{itemize}
    \item 若随机生成边时因达到最大尝试次数而未生成全部 \(E\) 条边，会输出警告信息。
\end{itemize}

\subsection{测试数据}
\begin{itemize}
    \item \textbf{正确输入}：例如 \(V=30, E=100\)，程序应生成图并输出平均值与最大值。
    \item \textbf{边界输入}：\(V=1, E=0\)（单个节点无边），平均值和最大值均为1，占比100%，应判定为耦合度高。
    \item \textbf{错误输入}：\(V=0, E=10\) 或 \(V=5, E=30\)（超过最大边数 \(5\times4=20\)），程序提示错误并退出。
\end{itemize}

\section{概要设计}

\subsection{问题解决思路}
程序采用邻接矩阵存储有向图，利用广度优先搜索计算每个顶点的可达节点数。整体流程如下：
\begin{enumerate}
    \item 接收用户输入的 \(V\) 和 \(E\)，进行合法性检查。
    \item 初始化 \(V \times V\) 的邻接矩阵，所有元素置0。
    \item 定义核心服务顶点集：前 \(10\%\) 的顶点（索引 \(0\) 至 \(coreCount-1\)），至少包含1个顶点。
    \item 随机生成 \(E\) 条有向边：
        \begin{itemize}
            \item 每条边以 \(50\%\) 的概率从核心服务中均匀选取源顶点，以 \(50\%\) 的概率从普通服务中选取。
            \item 目的顶点均匀随机选取。
            \item 禁止自环和重复边；若因冲突导致无法生成足够边数，会在达到最大尝试次数后停止并给出警告。
        \end{itemize}
    \item 对每个顶点执行 BFS，统计从该顶点出发能到达的顶点数量。
    \item 计算所有可达数的平均值与最大值，输出并给出风险评估结论。
\end{enumerate}

\subsection{数据结构定义}
\begin{itemize}
    \item \textbf{邻接矩阵}：使用二维指针 \texttt{int **adj}，大小为 \(V \times V\)，\texttt{adj[u][v]=1} 表示存在从 \(u\) 到 \(v\) 的有向边。
    \item \textbf{队列（BFS用）}：
\begin{lstlisting}
typedef struct {
    int *data;
    int front, rear, size, capacity;
} Queue;
\end{lstlisting}
    \item \textbf{辅助布尔数组}：在 BFS 中标记顶点是否已访问。
\end{itemize}

\subsection{主程序流程与模块层次}
\begin{verbatim}
main
├── 提示并读取 V, E
├── 合法性检查
├── 初始化随机种子
├── 分配邻接矩阵并清零
├── 定义核心服务顶点数 coreCount
├── 生成 E 条边（带重试）
│   ├── 选择源顶点（核心或普通，概率各半）
│   ├── 随机选择目的顶点
│   ├── 检查自环和重复边
│   └── 添加边并更新计数
├── 对每个顶点 i 调用 bfsReachable
│   ├── 创建访问标记数组
│   ├── 创建队列，将 i 入队
│   ├── 循环出队，遍历所有邻接顶点
│   └── 统计可达顶点数并返回
├── 计算平均值、最大值
├── 输出结果并给出风险评估
└── 释放邻接矩阵内存
\end{verbatim}

\section{详细设计}

\subsection{核心算法伪代码}

\subsubsection{随机图生成（含核心服务偏好）}
\begin{lstlisting}
coreCount = max(V/10, 1)
edgesGenerated = 0
attempts = 0
while edgesGenerated < E and attempts < MAX_ATTEMPTS:
    if random() % 100 < 50:
        src = random() % coreCount            // 核心服务
    else:
        if V > coreCount:
            src = coreCount + random() % (V - coreCount) // 普通服务
        else:
            continue
    dst = random() % V
    if src == dst or adj[src][dst] == 1:
        attempts++; continue
    adj[src][dst] = 1
    edgesGenerated++
    attempts = 0
if edgesGenerated < E:
    打印警告信息
\end{lstlisting}

\subsubsection{从顶点 s 出发的 BFS 可达计数}
\begin{lstlisting}
function bfsReachable(adj, V, s) -> int:
    visited = calloc(V, sizeof(bool))
    queue = createQueue(V)
    visited[s] = true
    enqueue(queue, s)
    count = 1
    while not isEmpty(queue):
        u = dequeue(queue)
        for v = 0 to V-1:
            if adj[u][v] == 1 and not visited[v]:
                visited[v] = true
                count++
                enqueue(queue, v)
    free(visited); freeQueue(queue)
    return count
\end{lstlisting}

\subsubsection{主控流程}
\begin{lstlisting}
输入 V, E
if V <= 0 或 E < 0 或 E > V*(V-1): 报错退出
adj = 分配 V×V 零矩阵
coreCount = V/10; if coreCount == 0: coreCount = 1
生成边（如上述）
sum = 0; maxReach = 0
for i in 0..V-1:
    reach = bfsReachable(adj, V, i)
    sum += reach
    if reach > maxReach: maxReach = reach
avg = sum / V
输出 avg, maxReach
if avg/V >= 0.5 or maxReach/V >= 0.5: 输出 "耦合度高"
else: 输出 "耦合度低"
释放 adj
\end{lstlisting}

\subsection{函数调用关系图}
\begin{verbatim}
main
├── createQueue / enqueue / dequeue / isEmpty / freeQueue
│   └── (在 bfsReachable 中调用)
├── bfsReachable
│   ├── calloc / free (visited)
│   └── 使用队列函数
└── 随机数生成 (rand / srand)
\end{verbatim}

\section{调试分析报告}

\subsection{调试过程中遇到的问题及解决方法}
\begin{enumerate}
    \item \textbf{问题1：重试机制导致死循环或效率低下} \\
    当 \(E\) 接近最大可能边数 \(V(V-1)\) 时，随机选择源和目的很容易遇到已存在的边，导致重试次数急剧增加。 \\
    \textbf{解决方法}：设置最大尝试次数 \(10^6\)，若达到上限则停止添加边并给出警告，同时保留已生成的边。这样可以避免无限循环，且程序仍可继续计算（虽然边数可能少于要求，但不影响演示目的）。

    \item \textbf{问题2：核心服务占比小于1时的处理} \\
    当 \(V < 10\) 时，\(V/10\) 为 0，但至少需要1个核心服务。 \\
    \textbf{解决方法}：在计算 \texttt{coreCount} 后判断，若为0则强制设为1。
\end{enumerate}

\subsection{算法时空复杂度分析}
\begin{itemize}
    \item \textbf{时间复杂度}：
        \begin{itemize}
            \item 生成边：常数复杂度
            \item BFS 计算：对每个顶点运行一次 BFS，每次 BFS 遍历所有顶点和边，复杂度 \(O(V \cdot (V+E))\)。对于中等规模图（如 \(V \leq 1000\)）可接受。
        \end{itemize}
    \item \textbf{空间复杂度}：
        \begin{itemize}
            \item 邻接矩阵：\(V^2\) 个整数，约为 \(4V^2\) 字节。
            \item BFS 辅助数组：\(O(V)\)。
            \item 队列：\(O(V)\)。
            \item 总空间复杂度 \(O(V^2)\)，当 \(V\) 较大时内存占用显著。
        \end{itemize}
\end{itemize}

\subsection{改进设想}
\begin{itemize}
    \item 对于大规模图，邻接矩阵内存过大，可以在边较少时改用邻接表存储。
    \item 增加对生成图的统计信息输出，如平均出度、入度分布等。
    \item 支持从文件读取边列表，而非随机生成，以分析真实拓扑。
\end{itemize}

\section{用户使用说明}

\begin{enumerate}
    \item 编译程序：\texttt{gcc -o graph\_exp chap7.exp.c -O2}
    \item 运行程序：\texttt{./graph\_exp}
    \item 根据提示输入顶点数 \(V\) 和边数 \(E\)，例如：
\begin{verbatim}
请输入顶点数 V（正整数，取值范围：1 ~ 任意正整数）: 30
请输入边数 E（非负整数，最大可能边数为 V*(V-1)）: 100
\end{verbatim}
    \item 程序输出平均值和最大值，并给出耦合性结论，例如：
\begin{verbatim}
平均值为 16.23，最大值为 28，说明系统耦合度高
\end{verbatim}
    \item 若边数超过最大可能值，程序会提示错误并退出。
\end{enumerate}

\section{测试结果}

\subsection{测试环境}
\begin{itemize}
    \item 操作系统：Windows 11
    \item 编译器：MinGW-W64 GCC 8.1.0 / GCC 9.4.0
    \item 编译选项：\texttt{-std=c11 -Wall}
\end{itemize}

\subsection{测试用例1：中等规模图（V=30, E=100）}
\textbf{测试步骤}：运行程序，输入 30 和 100。\\
\textbf{结果}：\texttt{平均值为 26.13，最大值为 30，说明系统耦合度高}

\subsection{测试用例2：稀疏图（V=50, E=30）}
\textbf{测试数据}：50 个顶点，30 条边。\\
\textbf{结果}：\texttt{平均值为 1.78，最大值为 14，说明系统耦合度低}

\subsection{测试用例3：完全图（V=10, E=90）}
\textbf{测试数据}：10 个顶点，最大边数为 90，E=90 表示完全有向图（无自环）。\\
\textbf{结果}：\texttt{平均值为 10.00，最大值为 10，说明系统耦合度低}

\subsection{测试用例4：单顶点（V=1, E=0）}
\textbf{测试数据}：1 个顶点，0 条边。\\
\textbf{结果}：\texttt{平均值为 1.00，最大值为 1，说明系统耦合度高}

\subsection{测试用例5：非法输入（V=0, E=5）}
\textbf{测试数据}：顶点数 0。\\
\textbf{结果}：程序输出“错误: V 必须为正整数（≥1）。”并退出。符合预期。

\subsection{测试用例6：边数超限（V=5, E=30）}
\textbf{测试数据}：5 个顶点，最大边数为 20，输入 30。\\
\textbf{结果}：程序输出“错误: E 不能超过 V*(V-1)=20。”并退出。错误检测有效。

\end{document}